\subsection{Surface Reconstruction from Point Clouds}\label{subsec:nksr}

Surface reconstruction addresses the inverse problem of recovering a
continuous, connected 3D surface from a discrete, unstructured set of
points sampled from that surface.
In practical settings such point clouds are often sparse, non-uniformly
sampled, and corrupted by noise, containing no connectivity information
and no guarantee of watertightness.
Converting such observations into a renderable or printable mesh is a
well-studied problem in geometry processing, and one where recent
learning-based methods have demonstrated substantial advantages over
classical approaches.

\subsubsection{Classical Baseline: Poisson Surface Reconstruction}

Poisson Surface Reconstruction (PSR) \cite{kazhdan2006poisson} frames
mesh generation as a global Poisson equation.
Given a point cloud augmented with per-point inward-facing normals
(estimated, e.g., via \acrshort{pca} over a local $k$-nearest-neighbour
neighbourhood), PSR solves for an implicit indicator function $\chi$
whose gradient best matches the input normals.
A mesh is then extracted at an isovalue γ chosen as the mean of χ
evaluated at the input sample positions.
The method produces smooth, globally consistent surfaces and is well
understood, but it is sensitive to normal estimation quality and
requires a globally consistent orientation of input normals.

\subsubsection{Neural Kernel Surface Reconstruction (NKSR)}

NKSR \cite{huang2023nksr} is a learning-based method that replaces the
classical \acrshort{pde} solve with a learned neural prior, building on
the Neural Kernel Fields (NKF) framework while addressing its two main
limitations: inability to scale beyond approximately ten thousand input
points, and poor robustness to noise.

Like PSR, NKSR requires an \emph{oriented} point cloud as input, meaning
per-point surface normals must be provided or estimated beforehand.
The reconstructed surface is encoded as the zero level set of a
\emph{Neural Kernel Field} --- an implicit function expressed as a
weighted sum of learned, positive-definite kernel basis functions
conditioned on the input and centered on voxels in a predicted sparse
hierarchy.
A sparse convolutional backbone processes the input and predicts a
multi-scale voxel hierarchy in which each voxel carries a feature vector
and a predicted normal.
These define a set of compactly supported kernel basis functions whose
scalar coefficients are found by solving a sparse, positive-definite
linear system.
The system minimises a loss that simultaneously encourages (i) the
field value to vanish at all input points, and (ii) the field's gradient
to agree with the predicted normals at voxel centres near the surface.
The second term --- absent in the original NKF --- is what confers
robustness to noise: rather than interpolating exact positional
constraints, the model fits a smooth gradient field, making it
significantly less sensitive to measurement errors.
Compact kernel support enables the use of memory-efficient sparse linear
solvers, allowing NKSR to reconstruct millions of points in seconds.
The final mesh is extracted via dual Marching Cubes, and an optional
learned masking branch discards spurious geometry at coarse hierarchy
levels far from the surface.

NKSR is trained directly from dense oriented point clouds and requires
neither watertight meshes nor pre-computed occupancy labels as
supervision, making it applicable to a wide variety of acquisition
settings.
The result is a high-quality mesh that accurately reflects the overall
shape suggested by the input, even under sparse or noisy conditions.
